\frontsection{Tóm tắt}

Già hóa dân số làm tăng nhu cầu hỗ trợ người cao tuổi tại nhà, trong khi người thân không phải lúc nào cũng có thể hiện diện trực tiếp. Đồ án nghiên cứu và xây dựng một robot trợ lý ảo sử dụng thiết bị Android gắn trên robot làm nền tảng cảm biến, tính toán và tương tác. Mục tiêu là kết hợp giám sát an toàn, nhắc lịch, hội thoại, liên lạc và điều khiển từ xa trong một nguyên mẫu có chi phí phần cứng hợp lý.

Hệ thống gồm RobotApp, CaregiverApp, khối mô hình AI và firmware ESP32. RobotApp xử lý camera tại biên bằng mô hình ước lượng tư thế TensorFlow Lite, theo dõi người bằng SORT/Kalman/Hungarian và tái định danh nhẹ, sau đó đưa chuỗi 15 khung hình với 69 đặc trưng vào LSTM. Bộ hậu kiểm sử dụng vùng bao, keypoint, tỷ lệ tư thế, vị trí thân trên và trạng thái mất dấu để quyết định có hỏi xác nhận hay không. Khi người dùng cần trợ giúp hoặc không phản hồi, ứng dụng ghi cảnh báo và khởi tạo cuộc gọi khẩn cấp. RobotApp còn tích hợp nhận dạng giọng nói, Groq Chat Completions, Android Text-to-Speech và dịch vụ nhắc lịch.

CaregiverApp cho phép quản lý hồ sơ, lịch nhắc, cờ theo dõi/phát hiện té ngã, nhật ký cảnh báo--nhắc việc--cuộc gọi, gọi video hai chiều và điều khiển robot bằng joystick. Firebase Authentication và Realtime Database được dùng cho danh tính, dữ liệu, lệnh và signaling; âm thanh/hình ảnh truyền bằng WebRTC. RobotApp kiểm tra timestamp và miền lệnh trước khi chuyển thành tốc độ hai bánh rồi gửi qua BLE GATT. Firmware ESP32 sử dụng NimBLE, SimpleFOC, hai encoder AS5600 và hai động cơ BLDC; watchdog 800 ms đưa robot về dừng khi mất lệnh.

Kết quả là nguyên mẫu tích hợp xuyên suốt từ giao diện người chăm sóc tới phần thân robot, kèm các artifact mô hình H5/PT/ONNX/TFLite và quy trình chuyển đổi cho Android. Báo cáo đặc tả các cơ chế an toàn, ma trận truy vết và phương pháp đánh giá; đồng thời công bố giới hạn về dữ liệu thực địa, metric AI, node Firebase toàn cục, phụ thuộc dịch vụ Internet và bảo mật trước khi phát triển thành sản phẩm.

\noindent\textbf{Từ khóa:} \textit{robot trợ lý, người cao tuổi, phát hiện té ngã, pose estimation, LSTM, Android, WebRTC, BLE, điều khiển FOC.}

\frontsection{Abstract}

Population ageing increases the demand for in-home assistance while family caregivers cannot always be physically present. This thesis investigates and develops a virtual-assistant robot in which an Android device mounted on the robot provides sensing, computation, communication, and interaction. The objective is to integrate safety monitoring, reminders, voice dialogue, live communication, and remote motion control into a cost-conscious prototype.

The system consists of RobotApp, CaregiverApp, an AI model toolchain, and ESP32 firmware. RobotApp performs on-device pose inference with TensorFlow Lite, tracks the primary person using SORT, Kalman filtering, Hungarian assignment, and lightweight re-identification, and feeds a 15-frame sequence of 69 features to an LSTM. A verification layer examines bounding-box integrity, keypoint visibility, body geometry, upper-body position, and target loss before requesting confirmation. If the person asks for help or does not respond, the system records an alert and starts an emergency call. RobotApp also integrates speech recognition, Groq Chat Completions, Android Text-to-Speech, and reminder services.

CaregiverApp supports profile and schedule management, tracking and fall-detection controls, reminder/alert/call diaries, two-way video calls, and a virtual joystick. Firebase Authentication and Realtime Database provide identity, shared data, commands, and signaling, while WebRTC transports audio and video. RobotApp validates command ranges and timestamps before mapping commands to two wheel targets and sending them over BLE GATT. The ESP32 firmware uses NimBLE, SimpleFOC, two AS5600 encoders, and two BLDC motors; an 800-ms watchdog stops the robot when commands are lost.

The outcome is an end-to-end prototype spanning caregiver interaction, edge AI, real-time services, and the physical robot body, together with H5/PT/ONNX/TFLite artifacts and a model-conversion workflow. The thesis specifies safety mechanisms, requirement traceability, and an evaluation protocol, while explicitly identifying limitations in field data, quantitative AI metrics, globally scoped Firebase nodes, Internet dependencies, and security hardening required before product deployment.

\noindent\textbf{Keywords:} \textit{assistive robot, older adults, fall detection, pose estimation, LSTM, Android, WebRTC, BLE, field-oriented control.}
